
    %%%%%%%%%%%%%%%%%%%%%%%%%%%%%%%%%%%%%%%%%%%%%%%%%%%%%%%%%%%%%%%%%%%%%%%%%%%%%%%%
    %%% Classe do documento
    \documentclass[12pt, addpoints]{exam}
    \UseRawInputEncoding

    %%%%%%%%%%%%%%%%%%%%%%%%%%%%%%%%%%%%%%%%%%%%%%%%%%%%%%%%%%%%%%%%%%%%%%%%%%%%%%%%
    % preambulo.tex
    %
    % Arquivo de configuração de packages para uso o LaTeX, conforme minhas
    % preferências e modelos pessoais.
    %
    % Para maiores informações, visite:
    %    https://github.com/abrantesasf/latex
    %
    % NÃO ALTERE SE NÃO SOUBER O QUE ESTÁ FAZENDO!
    %%%%%%%%%%%%%%%%%%%%%%%%%%%%%%%%%%%%%%%%%%%%%%%%%%%%%%%%%%%%%%%%%%%%%%%%%%%%%%%%


    %%%%%%%%%%%%%%%%%%%%%%%%%%%%%%%%%%%%%%%%%%%%%%%%%%%%%%%%%%%%%%%%%%%%%%%%%%%%%%%%
    %%% Estruturas de controle:
    \input{static/utils/controles}

    %%%%%%%%%%%%%%%%%%%%%%%%%%%%%%%%%%%%%%%%%%%%%%%%%%%%%%%%%%%%%%%%%%%%%%%%%%%%%%%%
    %%% Compilação condicional em PDF:
    \input{static/utils/pdf}

    %%%%%%%%%%%%%%%%%%%%%%%%%%%%%%%%%%%%%%%%%%%%%%%%%%%%%%%%%%%%%%%%%%%%%%%%%%%%%%%%
    %%% Configurações de layout da página:
    \input{static/utils/layout}

    %%%%%%%%%%%%%%%%%%%%%%%%%%%%%%%%%%%%%%%%%%%%%%%%%%%%%%%%%%%%%%%%%%%%%%%%%%%%%%%%
    %%% Configurações para autores:
    \input{static/utils/autores}

    %%%%%%%%%%%%%%%%%%%%%%%%%%%%%%%%%%%%%%%%%%%%%%%%%%%%%%%%%%%%%%%%%%%%%%%%%%%%%%%%
    %%% Cabeçalho e rodapé:
    %%% Atenção! É MUITO DIFÍCIL ajustar apropriadamente os running headers
    %%% e running footers da primeira vez. Você pode confiar no LaTeX e deixar que
    %%% ele faça o ajuste automaticamente ou pode quebrar a cabeça e
    %%% tentar melhorar algo que o LaTeX já faz muito bem, mesmo que não fique
    %%% exatamente do jeito que você quer. O que você prefere? Se quiser ajustar
    %%% manualmente, descomente a linha abaixo e faça as configurações no arquivo
    %%% "utils/headings.tex".
    %\input{static/utils/headings}

    %%%%%%%%%%%%%%%%%%%%%%%%%%%%%%%%%%%%%%%%%%%%%%%%%%%%%%%%%%%%%%%%%%%%%%%%%%%%%%%%
    %%% Configuração para as fontes do documento:
    %%% Se você quiser usar fontes próprias ou do sistema, precisará ajustar as
    %%% configurações no arquivo "utils/fontes.tex".
    %%% ATENÇÃO: por padrão eu utilizo XeLaTeX com fontes proprietárias projetadas
    %%% por Matthew Butterick, adquiridas em https://mbtype.com, que não podem ser
    %%% distribuídas devido à licença de utilização. Se você usar XeLaTeX, você
    %%% DEVERÁ OBRIGATORIAMENTE ajustar as fontes no arquivo "utils/fontes.tex".
    \input{static/utils/fontes}

    %%%%%%%%%%%%%%%%%%%%%%%%%%%%%%%%%%%%%%%%%%%%%%%%%%%%%%%%%%%%%%%%%%%%%%%%%%%%%%%%
    %%% Sumário:
    \input{static/utils/toc}

    %%%%%%%%%%%%%%%%%%%%%%%%%%%%%%%%%%%%%%%%%%%%%%%%%%%%%%%%%%%%%%%%%%%%%%%%%%%%%%%%
    %%% Matemática:
    \input{static/utils/matematica}

    %%%%%%%%%%%%%%%%%%%%%%%%%%%%%%%%%%%%%%%%%%%%%%%%%%%%%%%%%%%%%%%%%%%%%%%%%%%%%%%%
    %%% Notas de rodapé:
    \input{static/utils/footnotes}

    %%%%%%%%%%%%%%%%%%%%%%%%%%%%%%%%%%%%%%%%%%%%%%%%%%%%%%%%%%%%%%%%%%%%%%%%%%%%%%%%
    %%% Referências cruzadas, links, citações:
    \input{static/utils/crossrefs}

    %%%%%%%%%%%%%%%%%%%%%%%%%%%%%%%%%%%%%%%%%%%%%%%%%%%%%%%%%%%%%%%%%%%%%%%%%%%%%%%%
    %%% Configurações para os hiperlinks em PDF:
    \input{static/utils/pdfconfig}

    %%%%%%%%%%%%%%%%%%%%%%%%%%%%%%%%%%%%%%%%%%%%%%%%%%%%%%%%%%%%%%%%%%%%%%%%%%%%%%%%
    %%% Glossário e índice remissivo:
    \input{static/utils/glossindex}

    %%%%%%%%%%%%%%%%%%%%%%%%%%%%%%%%%%%%%%%%%%%%%%%%%%%%%%%%%%%%%%%%%%%%%%%%%%%%%%%%
    %%% Referências bibliográficas:
    \input{static/utils/bibliografia}

    %%%%%%%%%%%%%%%%%%%%%%%%%%%%%%%%%%%%%%%%%%%%%%%%%%%%%%%%%%%%%%%%%%%%%%%%%%%%%%%%
    %%% Cores:
    \input{static/utils/cores}

    %%%%%%%%%%%%%%%%%%%%%%%%%%%%%%%%%%%%%%%%%%%%%%%%%%%%%%%%%%%%%%%%%%%%%%%%%%%%%%%%
    %%% Computação:
    \input{static/utils/computacao}

    %%%%%%%%%%%%%%%%%%%%%%%%%%%%%%%%%%%%%%%%%%%%%%%%%%%%%%%%%%%%%%%%%%%%%%%%%%%%%%%%
    %%% Gráficos:
    \input{static/utils/graficos}

    %%%%%%%%%%%%%%%%%%%%%%%%%%%%%%%%%%%%%%%%%%%%%%%%%%%%%%%%%%%%%%%%%%%%%%%%%%%%%%%%
    %%% Ícones:
    \input{static/utils/icones}

    %%%%%%%%%%%%%%%%%%%%%%%%%%%%%%%%%%%%%%%%%%%%%%%%%%%%%%%%%%%%%%%%%%%%%%%%%%%%%%%%
    %%% Blocos:
    \input{static/utils/blocos}

    %%%%%%%%%%%%%%%%%%%%%%%%%%%%%%%%%%%%%%%%%%%%%%%%%%%%%%%%%%%%%%%%%%%%%%%%%%%%%%%%
    %%% Boxes coloridos:
    \input{static/utils/colorbox}

    %%%%%%%%%%%%%%%%%%%%%%%%%%%%%%%%%%%%%%%%%%%%%%%%%%%%%%%%%%%%%%%%%%%%%%%%%%%%%%%%
    %%% Tabelas:
    \input{static/utils/tabelas}

    %%%%%%%%%%%%%%%%%%%%%%%%%%%%%%%%%%%%%%%%%%%%%%%%%%%%%%%%%%%%%%%%%%%%%%%%%%%%%%%%
    %%% Ambientes de listas:
    \input{static/utils/listas}

    %%%%%%%%%%%%%%%%%%%%%%%%%%%%%%%%%%%%%%%%%%%%%%%%%%%%%%%%%%%%%%%%%%%%%%%%%%%%%%%%
    %%% Ambientes floats e similares:
    \input{static/utils/floats}

    %%%%%%%%%%%%%%%%%%%%%%%%%%%%%%%%%%%%%%%%%%%%%%%%%%%%%%%%%%%%%%%%%%%%%%%%%%%%%%%%
    %%% Ajustes para a classe EXAM:
    \input{static/utils/exam}

    %%%%%%%%%%%%%%%%%%%%%%%%%%%%%%%%%%%%%%%%%%%%%%%%%%%%%%%%%%%%%%%%%%%%%%%%%%%%%%%%
    %%% Ajustes para textos:
    \input{static/utils/texto}

    %%%%%%%%%%%%%%%%%%%%%%%%%%%%%%%%%%%%%%%%%%%%%%%%%%%%%%%%%%%%%%%%%%%%%%%%%%%%%%%%
    %%% Ambientes, aliases, comandos e customizações em geral para serem
    %%% utilizadas nos documentos. Defina aqui qualquer customização específica
    %%% para seus documentos!
    \input{static/utils/customizacoes}

    %%%%%%%%%%%%%%%%%%%%%%%%%%%%%%%%%%%%%%%%%%%%%%%%%%%%%%%%%%%%%%%%%%%%%%%%%%%%%%%%
    %%% Ajuste do layout, espaçamento de linhas e etc.:
    \geometry{a4paper, portrait, left=2.0cm, right=2.0cm, top=2.0cm, bottom=2.0cm}
    \singlespacing

    %%%%%%%%%%%%%%%%%%%%%%%%%%%%%%%%%%%%%%%%%%%%%%%%%%%%%%%%%%%%%%%%%%%%%%%%%%%%%%%%
    %%% Configurações de cabeçalho e rodapé (não use fancyhdr pois ocorre conflito):
    \pagestyle{headandfoot}
    \runningheadrule
    \firstpageheader{}{}{}
    \runningheader{Sem disciplina}{}{Avaliação}
    \firstpagefooter{}{}{Página \thepage\ de \numpages}
    \runningfooter{}{}{Página \thepage\ de \numpages}
    \runningfootrule

    %%%%%%%%%%%%%%%%%%%%%%%%%%%%%%%%%%%%%%%%%%%%%%%%%%%%%%%%%%%%%%%%%%%%%%%%%%%%%%%%
    %%% Configurações para as propriedades do PDF:
    \hypersetup{
    hidelinks,           % Comente para web, descomente para impressão
    colorlinks=true,      % True para web, False para impressão
    pdftitle=MARKINMACONHA: Avaliação,
    pdfauthor=Matheus Endlich Silveira,
    pdfsubject=Sem disciplina,
    pdfkeywords=['Sem tag'],
    pdfinfo={
        CreationDate={}, % Ex.: D:AAAAMMDDHH24MISS
        ModDate={}       % Ex.: D:AAAAMMDDHH24MISS
    }
    }
    \input{static/utils/pdfconfig}

    %%%%%%%%%%%%%%%%%%%%%%%%%%%%%%%%%%%%%%%%%%%%%%%%%%%%%%%%%%%%%%%%%%%%%%%%%%%%%%%%
    %%% Começa o documento
    \begin{document}

    %%%%%%%%%%%%%%%%%%%%%%%%%%%%%%%%%%%%%%%%%%%%%%%%%%%%%%%%%%%%%%%%%%%%%%%%%%%%%%%%
    %%% Folha de instruções padronizadas da UVV para a prova
    \pdfbookmark[1]{Instruções}{instrucoes}
    \begin{figure}[H]
        \begin{center}
            \includegraphics[scale=0.3]{{static/imgs/logo-uvv.png}}
        \end{center}	
    \end{figure}

    \vspace{-1cm}

    \begin{table}[H]
        \centering
        \begin{tabular}{|llll|}
            \hline
            \multicolumn{2}{|l|}{\textbf{Disciplina:} Sem disciplina} & 
            \multicolumn{1}{l|}{\multirow{3}{*}{\textbf{\makecell{Nota:\\ \,\\ \,}}}} & 
            \multirow{3}{*}{\textbf{\makecell{Coordenador:\\ \,\\ \,}}} \\ 
            \cline{1-2}
            \multicolumn{2}{|l|}{\textbf{Professor:} Matheus Endlich Silveira \hspace{4.72cm} } & 
            \multicolumn{1}{l|}{} & \\ 
            \cline{1-2}
            \multicolumn{2}{|l|}{\textbf{Aluno:}} & 
            \multicolumn{1}{l|}{} & \\ 
            \hline
            \multicolumn{1}{|l|}{\textbf{Turma:} \hspace{6.3cm} } & 
            \multicolumn{1}{l|}{\textbf{Semestre:}} & 
            \multicolumn{2}{l|}{\textbf{Valor:} 10 pontos} \\ 
            \hline
            \multicolumn{1}{|l|}{\textbf{Data:}} & 
            \multicolumn{3}{l|}{\textbf{Avaliação:}} \\ 
            \hline
        \end{tabular}
    \end{table}

    \begin{center}
        \fbox{\setlength\fboxsep{0.3cm} \fbox{\parbox{15.4cm}{
            \begin{center}
                \textbf{INSTRUÇÕES PARA A REALIZAÇÃO DESTA AVALIAÇÃO:}
            \end{center}
            \begin{itemize}[leftmargin=*]
                \item Esta é a primeira avaliação da disciplina Sem disciplina, 
                    e engloba o conteúdo referente Sem tag.
                \item Esta avaliação contém 15 questões objetivas, \textbf{todas obrigatórias}.
                \item \textbf{Leia atentamente cada questão!} As questões objetivas terão uma,
                    e apenas uma única, resposta a ser marcada.
                \item \textbf{Desligue o celular} antes de começar. A avaliação é
                    \textbf{individual} e \textbf{sem consulta}.
                \item As questões podem ser respondidas, na prova, com lápis ou caneta. Ao final
                    da prova \textbf{{VOCÊ DEVERÁ PREENCHER O GABARITO COM CANETA
                    DE COR PRETA}} \textbf{\underline{OU AZUL}} para que seja feita a correção
                    automatizada através da leitura do padrão de respostas marcado no
                    gabarito.
                \item \textbf{CUIDADO ao preencher o gabarito} pois eles são \textbf{INDIVIDUAIS
                    e NOMEADOS} (cada estudante tem um gabarito próprio específico, com um
                    código de identificação único). \textbf{Questões rasuradas são anuladas}
                    pelo software de leitura do gabarito.
                \item Ao final da prova, \textbf{DEVOLVA A PROVA E O GABARITO} para o professor.
                \item Siga todas as normas de \textbf{Integridade Acadêmica} da disciplina.
                    Alunos que forem flagrados com qualquer espécie de ``cola'' ou trocando
                    informações com outros alunos terão suas avaliações recolhidas, as notas
                    zeradas, e a situação será encaminhada para a coordenação para a aplicação
                    das penalidades previstas pela universidade.
                \item Com 90 minutos de avaliação você terá tempo de até 6,min por questão,
                    em média.
                \item Boa avaliação!
            \end{itemize}
            \vspace{2cm}
        }}}
    \end{center}
    \newpage

    %%%%%%%%%%%%%%%%%%%%%%%%%%%%%%%%%%%%%%%%%%%%%%%%%%%%%%%%%%%%%%%%%%%%%%%%%%%%%%%%
    %%% Inicia o bloco de questões:
    \begin{questions}

    \newpage
    \question

Qual dos protocolos abaixo pode ser caracterizado como protocolo de roteamento do tipo estado de enlace?
\begin{choices}
\choice RIP2
\choice ICMP
\choice IGMP
\choice OSPF
\choice BGP-4
\end{choices}

\question

(ENADE 2017) No modelo relacional de dados, uma junção entre tabelas cria uma

outra tabela derivada de duas ou mais tabelas de acordo com os critérios

especificados para a junção, de modo similar às regras da teoria dos

conjuntos. Nos conjuntos a seguir, as áreas hachuradas representam os resultados

das consultas SQL em que ``id' é uma chave primária e ``fk' é uma chave

estrangeira:

\begin{figure}[H]

  \centering

  \includegraphics[width=0.5\linewidth]{static/imgs/defaul.png}

\end{figure}

Assinale a opção em que a declaração SQL é equivalente ao conjunto apresentado a

seguir:

\begin{figure}[H]

  \centering

  \includegraphics[width=0.5\linewidth]{static/imgs/defaul.png}

\end{figure}


\begin{choices}
\choice \begin{verbatim}SELECT *

FROM tabela1 t1

WHERE NOT EXISTS (SELECT 1

                  FROM tabela1 t2

                  WHERE t2.id = tk.fk);

\end{verbatim}


\choice \begin{verbatim}SELECT *

FROM tabela1 t1

RIGHT JOIN tabela2 t2 ON (t1.id = t2.fk)

WHERE t1.id IS NULL;

\end{verbatim}


\choice \begin{verbatim}SELECT *

FROM tabela1 t1

RIGHT JOIN tabela2 t2 ON (t1.id = t2.fk)

WHERE t2.fk <> NULL;

\end{verbatim}


\choice \begin{verbatim}SELECT *

FROM tabela1 t1

INNER JOIN tabela2 t2 ON (t1.id = t2.fk);

\end{verbatim}


\choice \begin{verbatim}SELECT *

FROM tabela1 t1 WHERE EXISTS (SELECT 1

                              FROM tabela2 t2

                              WHERE t2.id = t1.fk);

\end{verbatim}


\end{choices}

\question

A arquitetura ilustrada na figura abaixo é um exemplo típico de qual sistema?

(CU = unidade de controle; IS = \ingles{stream} de instruções; PU = unidade de

processamento; DS = \ingles{stream} de dados; LM = memória local)

\begin{figure}[H]

\begin{center}

\includegraphics[width=0.5\linewidth]{static/imgs/defaul.png}

\end{center}

\end{figure}

\vspace{-1cm}
\begin{choices}
\choice Processadores multicore
\choice Multiprocessadores simétricos
\choice Acesso não uniforme à memória
\choice Clusters
\choice Processadores seqüenciais
\end{choices}

\question

(Adaptado do ENADE 2005) Avalie o seguite diagrama relacional de um banco de dados:

\begin{figure}[H]

  \centering

  \caption{Banco de dados de peças e fornecedores}

  \label{fig:pecas}

  %\fbox{

    \includegraphics[width=0.5\linewidth]{static/imgs/defaul.png}

  %}\\

  %\footnotesize{Fonte: xxxx.}

\end{figure}

Assinale a opção que apresenta um comando SQL que permite obter uma lista que

contenha o nome de cada fornecedor que tenha fornecido alguma peça, o código da

peça fornecida, a descrição dessa peça e a quantidade fornecida da referida

peça.
\begin{choices}
\choice \begin{verbatim}SELECT * FROM pecas INNER JOIN fornecedores INNER JOIN fornecimentos;\end{verbatim} 
\choice \begin{verbatim}SELECT * FROM pecas INNER JOIN fornecimentos f ON (p.codigo = f.cod_peca) INNER JOIN fornecedores f2 ON (f.cod_forn = f2.cod_forn) ORDER BY f2.nome;\end{verbatim} 
\choice \begin{verbatim}SELECT nome, codigo, descricao, quantidade FROM pecas INNER JOIN FORNECEDORES INNER JOIN FORNECIMENTOS;\end{verbatim} 
\choice \begin{verbatim}SELECT DISTINCT p.nome, p.codigo, p.descricao, f.quantidade FROM pecas p INNER JOIN fornecimentos f ON (p.codigo = f.cod_peca);\end{verbatim} 
\choice \begin{verbatim}SELECT f2.nome, p.codigo, p.descricao, f.quantidade FROM pecas INNER JOIN fornecimentos f ON (p.codigo = f.cod_peca) INNER JOIN fornecedores f2 ON (f.cod_forn = f2.cod_forn);\end{verbatim} 
\end{choices}

\question

Assinale o argumento válido, onde \( S_1 \), \( S_2 \) indicam premissas e \( S \) a conclusão:


\begin{choices}
\choice \( S_1 \): Se o cavalo estiver cansado então ele perderá a corrida \\

              \( S_2 \): O cavalo estava descansado \\

              \( S \): O cavalo ganhou a corrida
\choice \( S_1 \): Se o cavalo estiver cansado então ele perderá a corrida \\

              \( S_2 \): O cavalo estava descansado \\

              \( S \): O cavalo perdeu a corrida
\choice Se o cavalo estiver cansado então ele perderá a corrida \\

              \( S_2 \): O cavalo perdeu a corrida \\

              \( S \): O cavalo estava cansado
\choice \( S_1 \): Se o cavalo estiver cansado então ele perderá a corrida \\

              \( S_2 \): O cavalo ganhou a corrida \\

              \( S \): O cavalo estava descansado
\choice nenhuma das anteriores
\end{choices}

\question

Filtro da mediana é:
\begin{choices}
\choice Indicado para detectar formas específicas em imagens.
\choice Indicado para detectar bordas em imagens.
\choice Indicado para detectar tonalidades específicas em uma imagem.
\choice Nenhuma das respostas acima.
\choice Indicado para atenuar ruído com preservação de bordas (i.e., rápidas transições de nível em imagens).
\end{choices}

\question

Sobre a técnica conhecida como Z-buffer é correto afirmar que:
\begin{choices}
\choice As dimensões do Z-buffer são independentes das dimensões do frame buffer. 
\choice É possível realizar o cômputo das variáveis envolvidas de forma incremental.
\choice É uma técnica muito comum de detecção de colisão. 
\choice As primitivas geométricas precisam estar ordenadas de acordo com a distância em relação ao observador. 
\choice Nenhuma das alternativas acima está correta. 
\end{choices}

\question

Os Sistemas de Bancos de Dados (SBD) são essenciais hoje em dia, tanto para

pessoas físicas quanto para pessoas jurídicas. Assinale a alternativa que não

representa adequadamente a importância do uso de um SBD:
\begin{choices}
\choice Permitem que as empresas tenham mecanismos de obedecer à Lei Geral de Proteção de Dados Pessoais (LGPD), impedindo processos de anonimização das informações.
\choice Permitem que as empresas criem sites na Internet de modo que os próprios usuários consigam consultar e atualizar alguns dados.
\choice Permitem consultas históricas, de modo otimizado, para que alguma informação atual possa ser comparada com dados anteriores.
\choice ermitem que as empresas criem aplicativos para telefones celulares, oferecendo serviços aos seus usuários.
\choice Permitem às empresas coletar, armazenar, agregar, manipular e gerenciar dados, tomando assim melhores decisões de negócios e fornecendo diversos serviços úteis aos seus clientes.
\end{choices}

\question

São motivos importantes e gerais para o estudo das linguagens de programação,

exceto:
\begin{choices}
\choice Avanço geral da computação
\choice Melhorar o embasamento na escolha das linguagens para um projeto
\choice Entender a importância da implementação
\choice Aprender a sintaxe e a semântica de uma determinada linguagem
\choice Aumentar a capacidade de expressar idéias
\end{choices}

\question

Se \( \lim_{n \to \infty} \left( \frac{1}{n^2} + \frac{2}{n^2} + \cdots + \frac{n-1}{n^2} \right) = L \) então


\begin{choices}
\choice \( L = 0 \)
\choice \( L = \infty \)
\choice \( L = 1/2 \)
\choice \( L = 2 \)
\choice \( L = 1 \)
\end{choices}

\question

Para que serve a segmentação de um processador (\textit{pipelining})?
\begin{choices}
\choice simplificar a implementação do processador
\choice reduzir o número de instruções estáticas nos programas
\choice simplificar o conjunto de instruções
\choice permitir a execução de mais de uma instrução por ciclo de relógio
\choice aumentar a velocidade do relógio
\end{choices}

\question

Seja a seguinte linguagem, onde \(\epsilon\) representa o string vazio e \$ representa um marcador de fim de entrada:



$S \rightarrow ABCD \\$\\

$A \rightarrow a \mid \epsilon \\$\\

$B \rightarrow a \mid \epsilon \\$\\

$C \rightarrow c \mid \epsilon \\$\\

$D \rightarrow S \mid c \mid \epsilon$\\



É incorreto afirmar que:
\begin{choices}
\choice O conjunto FIRST(D) é igual ao conjunto FIRST(S)
\choice O conjunto FOLLOW(A) = a, c, \$
\choice O conjunto FOLLOW(D) é igual a FOLLOW(S)
\choice O conjunto FOLLOW(B) = c, \$
\choice O conjunto FIRST(A) = a, \(\epsilon\)
\end{choices}

\question

Em relação ao paradigma de programação cliente-servidor. Qual das afirmativas abaixo é FALSA?
\begin{choices}
\choice Um aplicativo servidor é um programa de propósito especial dedicado a fornecer um serviço, mas pode tratar de múltiplos clientes remotos ao mesmo tempo. 
\choice Um aplicativo cliente é um programa arbitrário que se torna temporariamente um cliente quando for necessário o acesso remoto a um serviço, mas também executa processamento local.
\choice Um aplicativo cliente pode acessar múltiplos serviços quando necessário.
\choice Um aplicativo servidor inicia ativamente o contato com clientes 
\choice Um aplicativo servidor aceita contato de clientes arbitrários, mas oferece um único serviço. 
\end{choices}

\question

Na linguagem SQL, qual das seguintes palavras-chave é usada para modificar os

dados em uma tabela existente?


\begin{choices}
\choice SELECT
\choice INSERT
\choice UPDATE
\choice Nenhuma das alternativas acima
\choice CREATE
\end{choices}

\question

Seja a árvore binária abaixo a representação de um espaço de estados para um problema \( p \), em que o estado inicial é \( a \), e \( i \) e \( f \) são estados finais.



\begin{figure}[H]

    \centering

    \includegraphics[width=0.5\linewidth]{static/imgs/defaul.png}

    \caption{Questão 59}

    \label{fig:enter-label}

\end{figure}



Um algoritmo de busca em largura-primeiro forneceria a seguinte sequência de estados como primeira alternativa a um caminho-solução para o problema \( p \):
\begin{choices}
\choice \( a \ c \ f \)
\choice \( a \ b \ c \ d \ e \ f \)
\choice \( a \ b \ d \ h \ e \ i \)
\choice \( a \ b \ d \ e \ f \)
\choice \( a \ b \ e \ i \)
\end{choices}



    %%%%%%%%%%%%%%%%%%%%%%%%%%%%%%%%%%%%%%%%%%%%%%%%%%%%%%%%%%%%%%%%%%%%%%%%%%%%%%%%
    %%% Finaliza o bloco de questões:
    \end{questions}
    \end{document}
    